\documentclass[class=article, crop=false, titlepage, twoside, multi={itemize, figure, verbatim}, float=false]{standalone}

\usepackage{import} % Required for importing other .tex docs.  (import uses everything bw Begin and End Doc)
\usepackage{float} % Required for specifying the exact location of a figure or table
\usepackage{graphicx} % Required for including images
\usepackage{wrapfig}
\usepackage[pdftex,breaklinks,colorlinks=true,linkcolor=black,citecolor=blue,urlcolor=red,linktocpage=false,pagebackref=true,filecolor=magenta]{hyperref}%http://www.tug.org/applications/hyperref/manual.html#x1-100003.6
\usepackage{cite}
\usepackage[toc,title,page]{appendix}
\usepackage{pdfpages} % enables loading a pdf into the doc
\usepackage{makeidx}
\usepackage{glossaries} % must be after hyperref
\usepackage{blindtext}
\usepackage{enumitem}
%\usepackage{caption}

%\setlist[description]{leftmargin=\parindent,labelindent=\parindent}

%\renewcommand*{\bibname}{References} % renames the bibliography

\newcommand{\HRule}{\rule{\linewidth}{0.5mm}} % Command to make the lines in the title page

\graphicspath{{img/}{GIS_ChampionSection/img/}{awardsChapter/GIS_ChampionSection/img/}{brandPart/awardsChapter/GIS_ChampionSection/img/}{img/}{pairedProgSection/img/}{methodChapter/pairedProgSection/img/}{methodPart/methodChapter/pairedProgSection/img/}{documentationSection/img/}{methodChapter/documentationSection/img/}{methodPart/methodChapter/documentationSection/img/}{docStorageOrgSection/img/}{methodChapter/docStorageOrgSection/img/}{methodPart/methodChapter/docStorageOrgSection/img/}{QGisSection/img/}{toolsChapter/QGisSection/img/}{servicePart/toolsChapter/QGisSection/img/}{ESRISection/img/}{toolChapter/ESRISection/img/}{servicePart/toolChapter/ESRISection/img/}{../../../../source/}{../../source/}{servicePart/applicationsChapter/treasurerSection/img/}}

%\setlength\parindent{0pt} % eliminates indents


\def\titlename{The hyperref Package\\ \medskip\large \LaTeX{} Packages }

\title{\HRule % Horizontal Line added
\\[.4cm] % space
\begin{figure}[H] % included image
\begin{center}	% centered horizontally
\includegraphics[scale=.45]{GIS_Logo_better.jpg}
\end{center}
\end{figure}
\Huge \bfseries \titlename \\ % Title text
\HRule \\[.4cm] % Horizontal Line added
\author{\Large Allegan County GIS \\\Large www.allegancounty.org/gis} % defines author
}  % inputs common title

\begin{document}% Document Begins

\ifstandalone
\maketitle % creates title page
\clearpage
\tableofcontents % creates TOC
\clearpage
\fi

\subsection[hyperref Package]{\LARGE hyperref Package}

\subsubsection[Introduction]{\Large Introduction}
\IfStandalone{\href{../../../../../documentation/packageDocs/hyperref2017.pdf}{Official hyperref package documentation}}{\href{../../documentation/packageDocs/hyperref2017.pdf}{Official hyperref package documentation}}

\textbf{Notes:}
\begin{itemize}
\item Add the \textit{hyperref package} to the preamble\begin{large}\textbf{ last}\end{large}~\cite{hyperref}
\item To use Tex in a pdf bookmark: use \begin{verbatim} \texorpdfstring{\\}{}

ie. \paragraph{Sample Text\texorpdfstring{\\}{}}\end{verbatim}
Creates a new line without an error.
\end{itemize}

\begin{verbatim}\usepackage[options]{hyperref}\end{verbatim}

\subsubsection[Simple Use]{\Large Simple Use}
\begin{verbatim}Use \href{URL}{DESCRIPTION} to add a link with description

\href{https://www.latex-tutorial.com}{Website with tutorials}
produces:
\end{verbatim}
\href{https://www.latex-tutorial.com}{Website with tutorials}

\subsubsection[Options]{\Large Options}
Add optional arguments to the usepackage line:\\
Useful options:\begin{itemize}
\item \begin{large}\textbf{pdftex}\end{large}\\enables other options like breaklines
\item \begin{large}\textbf{breaklinks}\end{large}\\allow links to be broken across several lines\\eg. \url{https://lists.gnu.org/archive/html/emacs-orgmode/2013-06/msg00776.html}
\item \begin{large}\textbf{colorlinks}\end{large}\\Colors the text of links and anchors.(default is false)
\item \begin{large}\textbf{linkcolor}\end{large}\\Color for normal internal links(default is red).
\item \begin{large}\textbf{anchorcolor}\end{large}\\Color for anchor text.
\item \begin{large}\textbf{citecolor}\end{large}\\Color for bibliographic citations in text.
\item \begin{large}\textbf{urlcolor}\end{large}\\Color for linked URLs

\end{itemize}

\subsubsection[Use with Options]{\Large Use with options}
\begin{verbatim}
\usepackage[breaklinks,colorlinks,citecolor=blue,
urlcolor=green]{hyperref}
\end{verbatim}

\subsubsection[Commands]{\Large Commands}
\begin{verbatim}
\href{URL}{text} — Makes text a link to URL.
\end{verbatim}

\begin{Large}To put a file path in text:\end{Large}\\
eg:

\IfStandalone{\href{../../../../../documentation/packageDocs/hyperref2017.pdf}{Official hyperref package documentation}}{\href{../../documentation/packageDocs/hyperref2017.pdf}{Official hyperref package documentation}}

\begin{tiny}(documentation Pt.4 pg.15)\end{tiny}
\begin{verbatim}
\href[options}{URL}{text}
\end{verbatim}

Options:\begin{itemize}
\item absolute
\begin{footnotesize}
\begin{verbatim}
\href{C:/AC/jalapeno/documentation/packageDocs/hyperref2017.pdf}
    {Official hyperref doc}
\end{verbatim}
\end{footnotesize}

\item relative \textbf{Note: relative path must be from final pdf location}
\begin{footnotesize}
\begin{verbatim}
\href{../../../../../documentation/packageDocs/hyperref2017.pdf}
    {Official hyperref package doc}
\end{verbatim}
This path works from main document
\begin{verbatim}

\href{../../documentation/packageDocs/hyperref2017.pdf}
    {Official hyperref package documentation}
\end{verbatim}
This path works from subsection document

\end{footnotesize}

\end{itemize}

\begin{verbatim}
\hyperref[label]{text}
    Makes text a link to where \ref{label} would point.

\hypertarget{name}{text}
    Sets an anchor on text with the label name.

\hyperlink{name}{text}
    Makes text a link that takes you to the anchor labeled name.
    Pair with \hypertarget.

\phantomsection
    Used in conjunction with

\addcontentsline
    to make the correct link in the Table of Contents.
\end{verbatim}

\end{document}
